% GENERATED FILE. Edit canonical lessons, then run npm run generate:books.
% canonical-source-hash: fnv1a64:97367f3a04ceb5e5
% canonical-lessons: HI-C342-qissa, HI-C342-atmakatha, HI-C342-prakasak, HI-C342-samskaran, HI-C342-anuvad, HI-R342-review-336-342, HI-R342-review-final-342

\chapter{Books and Their Makers}
\label{ch:hi-books-and-their-makers}
\hlchaptermodality{342}

\begin{chapteropening}
\textbf{By the end of this chapter:} \emph{I can talk about how a book came to be.}

Complete the last lesson of chapter 342: I can talk about how a book came to be.
\end{chapteropening}

\section[\texorpdfstring{\hi{क़िस्सा} (qissā)}{क़िस्सा (qissā)}]{\hi{क़िस्सा} (qissā) — an anecdote; a tale}

\label{lesson:HI-C342-qissa}



\begin{quote}
Before the new one: say the Hindi for a reader, then the Hindi for to be sold.
\end{quote}

\subsection*{You'll want to know: \hi{क़िस्सा}}
\textbf{\hi{क़िस्सा}} — \emph{qissā} — \textquotedblleft{}an anecdote; a tale\textquotedblright{}.

\begin{quote}
\hi{दादा} \hi{जी} \hi{के} \hi{पास} \hi{युद्ध} \hi{के} \hi{बहुत} \hi{क़िस्से} \hi{थे।}
\end{quote}

\begin{quote}
\emph{dādā jī ke pās yuddh ke bahut qisse the.}
\end{quote}

\begin{quote}
Grandfather had lots of stories from the war.
\end{quote}

Masculine.

\begin{cousinweb}
\textbf{\hi{क़िस्सा}} (\emph{qissā}) is usually traced, through Persian, to Arabic \emph{qiṣṣa}, \textquotedblleft{}a story\textquotedblright{}. It is a tale told aloud, often a true one from someone's life.
\end{cousinweb}

\begin{grammarlens}[title={What you've built}]
One more word for telling what happened.
\end{grammarlens}

\subsection*{Your turn}
\begin{itemize}
\raggedright
  \item \emph{Say it:} \emph{qissā}
  \item \emph{Say it:} \emph{qissā}, once more
  \item \emph{Say it:} \emph{dādā jī ke pās yuddh ke bahut qisse the.} — Grandfather had lots of stories from the war.
  \item \emph{Recall:} say the Hindi for a reader, then the Hindi for to be sold, then say \emph{qissā} again
\end{itemize}

\begin{tcolorbox}[breakable,colback=teal!4,colframe=teal!35!black,title={Before you move on}]
What does \hi{क़िस्सा} mean? (\textquotedblleft{}An anecdote; a tale\textquotedblright{}.) Say it once more.
\end{tcolorbox}
\section[\texorpdfstring{\hi{आत्मकथा} (ātmakathā)}{आत्मकथा (ātmakathā)}]{\hi{आत्मकथा} (ātmakathā) — an autobiography}

\label{lesson:HI-C342-atmakatha}



\begin{quote}
Before the new one: say the Hindi for to be sold, then the Hindi for an anecdote.
\end{quote}

\subsection*{You'll want to know: \hi{आत्मकथा}}
\textbf{\hi{आत्मकथा}} — \emph{ātmakathā} — \textquotedblleft{}an autobiography\textquotedblright{}.

\begin{quote}
\hi{उस} \hi{नेता} \hi{ने} \hi{अपनी} \hi{आत्मकथा} \hi{लिखी।}
\end{quote}

\begin{quote}
\emph{us netā ne apnī ātmakathā likhī.}
\end{quote}

\begin{quote}
That leader wrote his autobiography.
\end{quote}

Feminine.

\begin{cousinweb}
\textbf{\hi{आत्मकथा}} (\emph{ātmakathā}) joins Sanskrit \emph{ātman}, \textquotedblleft{}self\textquotedblright{}, and \textbf{\hi{कथा}}, a story (Monier-Williams): the story of oneself, told by oneself.
\end{cousinweb}

\begin{grammarlens}[title={What you've built}]
One more word for telling what happened.
\end{grammarlens}

\subsection*{Your turn}
\begin{itemize}
\raggedright
  \item \emph{Say it:} \emph{ātmakathā}
  \item \emph{Say it:} \emph{ātmakathā}, once more
  \item \emph{Say it:} \emph{us netā ne apnī ātmakathā likhī.} — That leader wrote his autobiography.
  \item \emph{Recall:} say the Hindi for to be sold, then the Hindi for an anecdote, then say \emph{ātmakathā} again
\end{itemize}

\begin{tcolorbox}[breakable,colback=teal!4,colframe=teal!35!black,title={Before you move on}]
What does \hi{आत्मकथा} mean? (\textquotedblleft{}An autobiography\textquotedblright{}.) Say it once more.
\end{tcolorbox}
\section[\texorpdfstring{\hi{प्रकाशक} (prakāśak)}{प्रकाशक (prakāśak)}]{\hi{प्रकाशक} (prakāśak) — a publisher}

\label{lesson:HI-C342-prakasak}



\begin{quote}
Before the new one: say the Hindi for an anecdote, then the Hindi for an autobiography.
\end{quote}

\subsection*{You'll want to know: \hi{प्रकाशक}}
\textbf{\hi{प्रकाशक}} — \emph{prakāśak} — \textquotedblleft{}a publisher\textquotedblright{}.

\begin{quote}
\hi{प्रकाशक} \hi{ने} \hi{वह} \hi{आत्मकथा} \hi{छापी।}
\end{quote}

\begin{quote}
\emph{prakāśak ne vah ātmakathā chāpī.}
\end{quote}

\begin{quote}
The publisher printed that autobiography.
\end{quote}

Masculine.

\begin{cousinweb}
\textbf{\hi{प्रकाशक}} (\emph{prakāśak}) is Sanskrit \emph{prakāśaka}, \textquotedblleft{}one who brings to light, who makes known\textquotedblright{} (Monier-Williams), from \emph{prakāśa}, light.
\end{cousinweb}

\begin{grammarlens}[title={What you've built}]
One more word for telling what happened.
\end{grammarlens}

\subsection*{Your turn}
\begin{itemize}
\raggedright
  \item \emph{Say it:} \emph{prakāśak}
  \item \emph{Say it:} \emph{prakāśak}, once more
  \item \emph{Say it:} \emph{prakāśak ne vah ātmakathā chāpī.} — The publisher printed that autobiography.
  \item \emph{Recall:} say the Hindi for an anecdote, then the Hindi for an autobiography, then say \emph{prakāśak} again
\end{itemize}

\begin{tcolorbox}[breakable,colback=teal!4,colframe=teal!35!black,title={Before you move on}]
What does \hi{प्रकाशक} mean? (\textquotedblleft{}A publisher\textquotedblright{}.) Say it once more.
\end{tcolorbox}
\section[\texorpdfstring{\hi{संस्करण} (saṁskaraṇ)}{संस्करण (saṁskaraṇ)}]{\hi{संस्करण} (saṁskaraṇ) — an edition}

\label{lesson:HI-C342-samskaran}



\begin{quote}
Before the new one: say the Hindi for an autobiography, then the Hindi for a publisher.
\end{quote}

\subsection*{You'll want to know: \hi{संस्करण}}
\textbf{\hi{संस्करण}} — \emph{saṁskaraṇ} — \textquotedblleft{}an edition\textquotedblright{}.

\begin{quote}
\hi{मुझे} \hi{उम्मीद} \hi{है} \hi{कि} \hi{इस} \hi{किताब} \hi{का} \hi{नया} \hi{संस्करण} \hi{अगले} \hi{साल} \hi{आएगा।}
\end{quote}

\begin{quote}
\emph{mujhe ummīd hai ki is kitāb kā nayā saṁskaraṇ agle sāl āegā.}
\end{quote}

\begin{quote}
I hope a new edition of this book will come out next year.
\end{quote}

Masculine.

\begin{cousinweb}
\textbf{\hi{संस्करण}} (\emph{saṁskaraṇ}) is Sanskrit \emph{saṁskaraṇa}, \textquotedblleft{}a preparing, a making ready\textquotedblright{} (Monier-Williams): a book made ready again.
\end{cousinweb}

\begin{grammarlens}[title={What you've built}]
One more word for telling what happened.
\end{grammarlens}

\subsection*{Your turn}
\begin{itemize}
\raggedright
  \item \emph{Say it:} \emph{saṁskaraṇ}
  \item \emph{Say it:} \emph{saṁskaraṇ}, once more
  \item \emph{Say it:} \emph{mujhe ummīd hai ki is kitāb kā nayā saṁskaraṇ agle sāl āegā.} — I hope a new edition of this book will come out next year.
  \item \emph{Recall:} say the Hindi for an autobiography, then the Hindi for a publisher, then say \emph{saṁskaraṇ} again
\end{itemize}

\begin{tcolorbox}[breakable,colback=teal!4,colframe=teal!35!black,title={Before you move on}]
What does \hi{संस्करण} mean? (\textquotedblleft{}An edition\textquotedblright{}.) Say it once more.
\end{tcolorbox}
\section[\texorpdfstring{\hi{अनुवाद} (anuvād)}{अनुवाद (anuvād)}]{\hi{अनुवाद} (anuvād) — a translation}

\label{lesson:HI-C342-anuvad}



\begin{quote}
Before the new one: say the Hindi for a publisher, then the Hindi for an edition.
\end{quote}

\subsection*{You'll want to know: \hi{अनुवाद}}
\textbf{\hi{अनुवाद}} — \emph{anuvād} — \textquotedblleft{}a translation\textquotedblright{}.

\begin{quote}
\hi{मीरा} \hi{ने} \hi{उस} \hi{कहानी} \hi{का} \hi{अंग्रेज़ी} \hi{में} \hi{अनुवाद} \hi{किया।}
\end{quote}

\begin{quote}
\emph{Mīrā ne us kahānī kā angrezī meṁ anuvād kiyā.}
\end{quote}

\begin{quote}
Mira translated that story into English.
\end{quote}

Masculine.

\begin{cousinweb}
\textbf{\hi{अनुवाद}} (\emph{anuvād}) is Sanskrit \emph{anuvāda}, \textquotedblleft{}a saying after, a repeating\textquotedblright{} (Monier-Williams): saying again in another language what was first said in one.
\end{cousinweb}

\begin{grammarlens}[title={What you've built}]
One more word for telling what happened.
\end{grammarlens}

\subsection*{Your turn}
\begin{itemize}
\raggedright
  \item \emph{Say it:} \emph{anuvād}
  \item \emph{Say it:} \emph{anuvād}, once more
  \item \emph{Say it:} \emph{Mīrā ne us kahānī kā angrezī meṁ anuvād kiyā.} — Mira translated that story into English.
  \item \emph{Recall:} say the Hindi for a publisher, then the Hindi for an edition, then say \emph{anuvād} again
\end{itemize}

\begin{tcolorbox}[breakable,colback=teal!4,colframe=teal!35!black,title={Before you move on}]
What does \hi{अनुवाद} mean? (\textquotedblleft{}A translation\textquotedblright{}.) Say it once more.
\end{tcolorbox}
\section[(dialogue)]{Looking back: from The Curse to Books and Their Makers}

\label{lesson:HI-R342-review-336-342}



\begin{quote}
Say the Hindi for penance and a curse.
\end{quote}

\subsection*{Your turn}
Say these aloud:

\begin{itemize}
\raggedright
  \item penance, a curse, to curse, to regret, to sulk
  \item a treasure, to bury, a cobra, a jewel, to bite (of a snake)
  \item a fairy, magic, a ghost, a witch, a miracle
  \item a prince, a disguise, to test, trickery, to lure
  \item a novel, a hero (of a story), a character (in a story), a chapter of a book, an outcome
  \item a page, a binding, a title, a reader, to be sold
  \item an anecdote, an autobiography, a publisher, an edition, a translation
\end{itemize}

\begin{tcolorbox}[breakable,colback=teal!4,colframe=teal!35!black,title={Before you move on}]
Penance? (\textbf{\hi{तपस्या}}, \emph{tapasyā}.) A curse? (\textbf{\hi{शाप}}, \emph{śāp}.) To curse? (\textbf{\hi{कोसना}}, \emph{kosnā}.) To regret? (\textbf{\hi{पछताना}}, \emph{pachtānā}.) To sulk? (\textbf{\hi{रूठना}}, \emph{rūṭhnā}.) A treasure? (\textbf{\hi{ख़ज़ाना}}, \emph{khazānā}.) To bury? (\textbf{\hi{गाड़ना}}, \emph{gāṛnā}.) A cobra? (\textbf{\hi{नाग}}, \emph{nāg}.) A jewel? (\textbf{\hi{मणि}}, \emph{maṇi}.) To bite (of a snake)? (\textbf{\hi{डसना}}, \emph{ḍasnā}.) A fairy? (\textbf{\hi{परी}}, \emph{parī}.) Magic? (\textbf{\hi{जादू}}, \emph{jādū}.) A ghost? (\textbf{\hi{भूत}}, \emph{bhūt}.) A witch? (\textbf{\hi{डायन}}, \emph{ḍāyan}.) A miracle? (\textbf{\hi{चमत्कार}}, \emph{camatkār}.) A prince? (\textbf{\hi{राजकुमार}}, \emph{rājkumār}.) A disguise? (\textbf{\hi{भेस}}, \emph{bhes}.) To test? (\textbf{\hi{परखना}}, \emph{parakhnā}.) Trickery? (\textbf{\hi{छल}}, \emph{chal}.) To lure? (\textbf{\hi{लुभाना}}, \emph{lubhānā}.) A novel? (\textbf{\hi{उपन्यास}}, \emph{upanyās}.) A hero (of a story)? (\textbf{\hi{नायक}}, \emph{nāyak}.) A character (in a story)? (\textbf{\hi{पात्र}}, \emph{pātra}.) A chapter of a book? (\textbf{\hi{अध्याय}}, \emph{adhyāy}.) An outcome? (\textbf{\hi{अंजाम}}, \emph{añjām}.) A page? (\textbf{\hi{पन्ना}}, \emph{pannā}.) A binding? (\textbf{\hi{जिल्द}}, \emph{jild}.) A title? (\textbf{\hi{शीर्षक}}, \emph{śīrṣak}.) A reader? (\textbf{\hi{पाठक}}, \emph{pāṭhak}.) To be sold? (\textbf{\hi{बिकना}}, \emph{biknā}.) An anecdote? (\textbf{\hi{क़िस्सा}}, \emph{qissā}.) An autobiography? (\textbf{\hi{आत्मकथा}}, \emph{ātmakathā}.) A publisher? (\textbf{\hi{प्रकाशक}}, \emph{prakāśak}.) An edition? (\textbf{\hi{संस्करण}}, \emph{saṁskaraṇ}.) A translation? (\textbf{\hi{अनुवाद}}, \emph{anuvād}.)
\end{tcolorbox}
\section[(dialogue)]{Books and Their Makers — once more}

\label{lesson:HI-R342-review-final-342}



\begin{quote}
Say the Hindi for an anecdote and an autobiography.
\end{quote}

\subsection*{Your turn}
\begin{itemize}
\raggedright
  \item \emph{Say it:} an anecdote, an autobiography, a publisher, an edition, a translation
\end{itemize}

\begin{tcolorbox}[breakable,colback=teal!4,colframe=teal!35!black,title={Before you move on}]
An anecdote? (\textbf{\hi{क़िस्सा}}, \emph{qissā}.) An autobiography? (\textbf{\hi{आत्मकथा}}, \emph{ātmakathā}.) A publisher? (\textbf{\hi{प्रकाशक}}, \emph{prakāśak}.) An edition? (\textbf{\hi{संस्करण}}, \emph{saṁskaraṇ}.) A translation? (\textbf{\hi{अनुवाद}}, \emph{anuvād}.)
\end{tcolorbox}
